% GENERATED FILE. Edit canonical lessons, then run npm run generate:books.
% canonical-source-hash: fnv1a64:71d1e22b783108b1
% canonical-lessons: LA-C469-incus, LA-C469-follis, LA-C469-caminus, LA-C469-ferrarius, LA-C469-cudo

\chapter{Anvil, Bellows, Forge, Blacksmith, To hammer}
\label{ch:la-anvil-bellows-forge-blacksmith-to-hammer}
\hlchaptermodality{469}

\begin{chapteropening}
\textbf{By the end of this chapter:} \emph{I can describe a smithy: the anvil, the bellows, the forge, the blacksmith and hammering iron.}

Complete the last lesson of chapter 469: I can describe a smithy: the anvil, the bellows, the forge, the blacksmith and hammering iron.
\end{chapteropening}

\section[incūs, incūdis]{incūs, incūdis — an anvil}

\label{lesson:LA-C469-incus}



\begin{quote}
Before the new one: say the Latin for a paint, then the Latin for a paintbrush.
\end{quote}

\subsection*{You'll want to know: incūs, incūdis}
\textbf{incūs, incūdis} — \textquotedblleft{}an anvil\textquotedblright{}.

\begin{quote}
Faber ferrum calidum in incūde pōnit.
\end{quote}

\begin{quote}
\emph{The smith puts the hot iron on the anvil.}
\end{quote}

Third declension, feminine: \emph{incūs, incūdis f.} \emph{Ad incūdem redīre}, to go back to the anvil, to rework a piece of writing. English \emph{incus} is a bone of the ear, named for its shape.

\begin{cousinweb}
\emph{incūs} is built on \emph{incūdō} 'to forge on' (EtymDB; Lewis \& Short), from \emph{in} 'on' and \emph{cūdō} 'to strike, to hammer': the block that iron is struck on.
\end{cousinweb}

\begin{grammarlens}[title={What you've built}]
One more word for telling, explaining and describing at B1.
\end{grammarlens}

\subsection*{Your turn}
\begin{itemize}
\raggedright
  \item \emph{Say it:} \emph{incūs, incūdis}
  \item \emph{Say it:} \emph{incūs, incūdis}, once more
  \item \emph{Say it:} a sentence of your own with \emph{incūs, incūdis}
  \item \emph{Recall:} say the Latin for a paint, then the Latin for a paintbrush, then say \emph{incūs, incūdis} again
\end{itemize}

\begin{tcolorbox}[breakable,colback=teal!4,colframe=teal!35!black,title={Before you move on}]
What does incūs, incūdis mean? (\textquotedblleft{}An anvil\textquotedblright{}.) Say it once more.
\end{tcolorbox}
\section[follis, follis]{follis, follis — a pair of bellows; a leather bag}

\label{lesson:LA-C469-follis}



\begin{quote}
Before the new one: say the Latin for a paintbrush, then the Latin for an anvil.
\end{quote}

\subsection*{You'll want to know: follis, follis}
\textbf{follis, follis} — \textquotedblleft{}a pair of bellows; a leather bag\textquotedblright{}.

\begin{quote}
Servus folle ignem auget.
\end{quote}

\begin{quote}
\emph{The slave makes the fire bigger with the bellows.}
\end{quote}

Third declension, masculine: \emph{follis, follis m.} Also a blown-up leather ball and a money bag. English \emph{fool} comes from it through French, from 'windbag'.

\begin{cousinweb}
\emph{follis} is inherited from PIE \emph{b\textsuperscript{h}olǵ\textsuperscript{h}nis} (EtymDB); Lewis \& Short compare \emph{flō} 'to blow'. English \emph{fool} comes from it through Old French \emph{fol}, from the sense 'windbag'.
\end{cousinweb}

\begin{grammarlens}[title={What you've built}]
One more word for telling, explaining and describing at B1.
\end{grammarlens}

\subsection*{Your turn}
\begin{itemize}
\raggedright
  \item \emph{Say it:} \emph{follis, follis}
  \item \emph{Say it:} \emph{follis, follis}, once more
  \item \emph{Say it:} a sentence of your own with \emph{follis, follis}
  \item \emph{Recall:} say the Latin for a paintbrush, then the Latin for an anvil, then say \emph{follis, follis} again
\end{itemize}

\begin{tcolorbox}[breakable,colback=teal!4,colframe=teal!35!black,title={Before you move on}]
What does follis, follis mean? (\textquotedblleft{}A pair of bellows; a leather bag\textquotedblright{}.) Say it once more.
\end{tcolorbox}
\section[camīnus, camīnī]{camīnus, camīnī — a forge, a furnace; a fireplace}

\label{lesson:LA-C469-caminus}



\begin{quote}
Before the new one: say the Latin for an anvil, then the Latin for a pair of bellows.
\end{quote}

\subsection*{You'll want to know: camīnus, camīnī}
\textbf{camīnus, camīnī} — \textquotedblleft{}a forge, a furnace; a fireplace\textquotedblright{}.

\begin{quote}
In camīnō ignis magnus ardet.
\end{quote}

\begin{quote}
\emph{A big fire burns in the forge.}
\end{quote}

Second declension, masculine: \emph{camīnus, camīnī m.} A smith's forge, or a stove for heating a room. English \emph{chimney} comes from a word built on it.

\begin{cousinweb}
\emph{camīnus} was borrowed from Greek \emph{kamīnos} 'furnace' (EtymDB; Lewis \& Short). English \emph{chimney} comes, through French, from Late Latin \emph{camīnāta}, built on it.
\end{cousinweb}

\begin{grammarlens}[title={What you've built}]
One more word for telling, explaining and describing at B1.
\end{grammarlens}

\subsection*{Your turn}
\begin{itemize}
\raggedright
  \item \emph{Say it:} \emph{camīnus, camīnī}
  \item \emph{Say it:} \emph{camīnus, camīnī}, once more
  \item \emph{Say it:} a sentence of your own with \emph{camīnus, camīnī}
  \item \emph{Recall:} say the Latin for an anvil, then the Latin for a pair of bellows, then say \emph{camīnus, camīnī} again
\end{itemize}

\begin{tcolorbox}[breakable,colback=teal!4,colframe=teal!35!black,title={Before you move on}]
What does camīnus, camīnī mean? (\textquotedblleft{}A forge, a furnace; a fireplace\textquotedblright{}.) Say it once more.
\end{tcolorbox}
\section[ferrārius, ferrāriī]{ferrārius, ferrāriī — a blacksmith, an ironworker}

\label{lesson:LA-C469-ferrarius}



\begin{quote}
Before the new one: say the Latin for a pair of bellows, then the Latin for a forge.
\end{quote}

\subsection*{You'll want to know: ferrārius, ferrāriī}
\textbf{ferrārius, ferrāriī} — \textquotedblleft{}a blacksmith, an ironworker\textquotedblright{}.

\begin{quote}
Ferrārius soleās ferreās equīs facit.
\end{quote}

\begin{quote}
\emph{The blacksmith makes iron shoes for the horses.}
\end{quote}

Second declension, masculine: \emph{ferrārius, ferrāriī m.}, short for \emph{faber ferrārius}. English \emph{farrier} comes from it through French.

\begin{cousinweb}
\emph{ferrārius} is built on \emph{ferrum} 'iron' with the suffix \emph{-ārius} (EtymDB; Lewis \& Short). English \emph{farrier}, a smith who shoes horses, comes from it through French.
\end{cousinweb}

\begin{grammarlens}[title={What you've built}]
One more word for telling, explaining and describing at B1.
\end{grammarlens}

\subsection*{Your turn}
\begin{itemize}
\raggedright
  \item \emph{Say it:} \emph{ferrārius, ferrāriī}
  \item \emph{Say it:} \emph{ferrārius, ferrāriī}, once more
  \item \emph{Say it:} a sentence of your own with \emph{ferrārius, ferrāriī}
  \item \emph{Recall:} say the Latin for a pair of bellows, then the Latin for a forge, then say \emph{ferrārius, ferrāriī} again
\end{itemize}

\begin{tcolorbox}[breakable,colback=teal!4,colframe=teal!35!black,title={Before you move on}]
What does ferrārius, ferrāriī mean? (\textquotedblleft{}A blacksmith, an ironworker\textquotedblright{}.) Say it once more.
\end{tcolorbox}
\section[cūdō, cūdere]{cūdō, cūdere — to hammer, to forge; to strike (coins)}

\label{lesson:LA-C469-cudo}



\begin{quote}
Before the new one: say the Latin for a forge, then the Latin for a blacksmith.
\end{quote}

\subsection*{You'll want to know: cūdō, cūdere}
\textbf{cūdō, cūdere} — \textquotedblleft{}to hammer, to forge; to strike (coins)\textquotedblright{}.

\begin{quote}
Faber ferrum in incūde cūdit.
\end{quote}

\begin{quote}
\emph{The smith hammers the iron on the anvil.}
\end{quote}

\emph{cūdō, cūdere, cūdī, cūsum}. \emph{Nummōs cūdere}, to strike coins.

\begin{cousinweb}
\emph{cūdō} is inherited from Proto-Italic \emph{kūdō} 'to stamp, to pound', from PIE \emph{kewh\textsubscript{2}-} 'to strike, to hew' (EtymDB). English \emph{hew} is a distant Germanic cousin.
\end{cousinweb}

\begin{grammarlens}[title={What you've built}]
One more word for telling, explaining and describing at B1.
\end{grammarlens}

\subsection*{Your turn}
\begin{itemize}
\raggedright
  \item \emph{Say it:} \emph{cūdō, cūdere}
  \item \emph{Say it:} \emph{cūdō, cūdere}, once more
  \item \emph{Say it:} a sentence of your own with \emph{cūdō, cūdere}
  \item \emph{Recall:} say the Latin for a forge, then the Latin for a blacksmith, then say \emph{cūdō, cūdere} again
\end{itemize}

\begin{tcolorbox}[breakable,colback=teal!4,colframe=teal!35!black,title={Before you move on}]
What does cūdō, cūdere mean? (\textquotedblleft{}To hammer, to forge; to strike (coins)\textquotedblright{}.) Say it once more.
\end{tcolorbox}
